% GENERATED FILE. Edit canonical lessons, then run npm run generate:books.
% canonical-source-hash: fnv1a64:d7cadc384802fd92
% canonical-lessons: JA-C277-ouyuki, JA-C277-tsumoru, JA-C277-fubuki, JA-C277-teiden, JA-C277-kogoeru

\chapter{Snowed In}
\label{ch:ja-snowed-in}
\hlchaptermodality{277}

\begin{chapteropening}
\textbf{By the end of this chapter:} \emph{I can say how deep the snow was and what it did.}

Complete the last lesson of chapter 277: I can say how deep the snow was and what it did.
\end{chapteropening}

\section[\texorpdfstring{\ja{おおゆき} (ōyuki)}{おおゆき (ōyuki)}]{\ja{おおゆき} (ōyuki) — heavy snow}

\label{lesson:JA-C277-ouyuki}



\begin{quote}
Before the new one: say the Japanese for a cliff, then the Japanese for awful.
\end{quote}

\subsection*{You'll want to know: \ja{おおゆき}}
\textbf{\ja{おおゆき}} — \emph{ōyuki} — \textquotedblleft{}heavy snow\textquotedblright{}.

\begin{quote}
\textbf{\ja{ことしは} \ja{おおゆきでした。}} — \emph{kotoshi wa ōyuki deshita} — \textquotedblleft{}We had a lot of snow this year.\textquotedblright{}
\end{quote}

\begin{cousinweb}
\textbf{\ja{おおゆき}} (\emph{ōyuki}) is a native Japanese word (wago), inherited rather than borrowed from Chinese, built from \emph{ō}, \textquotedblleft{}big\textquotedblright{}, and \emph{yuki}, \textquotedblleft{}snow\textquotedblright{} (JMdict, KANJIDIC2).
\end{cousinweb}

\begin{grammarlens}[title={What you've built}]
One more word for telling what happened.
\end{grammarlens}

\subsection*{Your turn}
\begin{itemize}
\raggedright
  \item \emph{Say it:} \emph{ōyuki}
  \item \emph{Say it:} \emph{ōyuki}, once more
  \item \emph{Say it:} \emph{kotoshi wa ōyuki deshita} — We had a lot of snow this year.
  \item \emph{Recall:} say the Japanese for a cliff, then the Japanese for awful, then say \emph{ōyuki} again
\end{itemize}

\begin{tcolorbox}[breakable,colback=teal!4,colframe=teal!35!black,title={Before you move on}]
What does \ja{おおゆき} mean? (\textquotedblleft{}Heavy snow\textquotedblright{}.) Say it once more.
\end{tcolorbox}
\section[\texorpdfstring{\ja{つもる} (tsumoru)}{つもる (tsumoru)}]{\ja{つもる} (tsumoru) — to accumulate; to lie (of snow)}

\label{lesson:JA-C277-tsumoru}



\begin{quote}
Before the new one: say the Japanese for awful, then the Japanese for heavy snow.
\end{quote}

\subsection*{You'll want to know: \ja{つもる}}
\textbf{\ja{つもる}} — \emph{tsumoru} — \textquotedblleft{}to accumulate; to lie (of snow)\textquotedblright{}.

\begin{quote}
\textbf{\ja{やねに} \ja{ゆきが} \ja{つもっています。}} — \emph{yane ni yuki ga tsumotte imasu} — \textquotedblleft{}There is snow lying on the roof.\textquotedblright{}
\end{quote}

Polite \textbf{\ja{つもります}}, \emph{tsumorimasu}; a godan verb. Past (polite): \textbf{\ja{ゆきが} \ja{つもりました。}}, \emph{yuki ga tsumorimashita}.

\begin{cousinweb}
\textbf{\ja{つもる}} (\emph{tsumoru}) is a native Japanese word (wago), inherited rather than borrowed from Chinese; its kanji was attached later as a kun'yomi (JMdict).
\end{cousinweb}

\begin{grammarlens}[title={What you've built}]
One more word for telling what happened.
\end{grammarlens}

\subsection*{Your turn}
\begin{itemize}
\raggedright
  \item \emph{Say it:} \emph{tsumoru}
  \item \emph{Say it:} \emph{tsumoru}, once more
  \item \emph{Say it:} \emph{yane ni yuki ga tsumotte imasu} — There is snow lying on the roof.
  \item \emph{Recall:} say the Japanese for awful, then the Japanese for heavy snow, then say \emph{tsumoru} again
\end{itemize}

\begin{tcolorbox}[breakable,colback=teal!4,colframe=teal!35!black,title={Before you move on}]
What does \ja{つもる} mean? (\textquotedblleft{}To accumulate; to lie (of snow)\textquotedblright{}.) Say it once more.
\end{tcolorbox}
\section[\texorpdfstring{\ja{ふぶき} (fubuki)}{ふぶき (fubuki)}]{\ja{ふぶき} (fubuki) — a snowstorm; a blizzard}

\label{lesson:JA-C277-fubuki}



\begin{quote}
Before the new one: say the Japanese for heavy snow, then the Japanese for to accumulate.
\end{quote}

\subsection*{You'll want to know: \ja{ふぶき}}
\textbf{\ja{ふぶき}} — \emph{fubuki} — \textquotedblleft{}a snowstorm; a blizzard\textquotedblright{}.

\begin{quote}
\textbf{\ja{ゆうべは} \ja{ひどい} \ja{ふぶきでした。}} — \emph{yūbe wa hidoi fubuki deshita} — \textquotedblleft{}There was a terrible blizzard last night.\textquotedblright{}
\end{quote}

\begin{cousinweb}
\textbf{\ja{ふぶき}} (\emph{fubuki}) is a native Japanese word (wago), inherited rather than borrowed from Chinese; its kanji was attached later as a kun'yomi (JMdict).
\end{cousinweb}

\begin{grammarlens}[title={What you've built}]
One more word for telling what happened.
\end{grammarlens}

\subsection*{Your turn}
\begin{itemize}
\raggedright
  \item \emph{Say it:} \emph{fubuki}
  \item \emph{Say it:} \emph{fubuki}, once more
  \item \emph{Say it:} \emph{yūbe wa hidoi fubuki deshita} — There was a terrible blizzard last night.
  \item \emph{Recall:} say the Japanese for heavy snow, then the Japanese for to accumulate, then say \emph{fubuki} again
\end{itemize}

\begin{tcolorbox}[breakable,colback=teal!4,colframe=teal!35!black,title={Before you move on}]
What does \ja{ふぶき} mean? (\textquotedblleft{}A snowstorm; a blizzard\textquotedblright{}.) Say it once more.
\end{tcolorbox}
\section[\texorpdfstring{\ja{ていでん} (teiden)}{ていでん (teiden)}]{\ja{ていでん} (teiden) — a power cut}

\label{lesson:JA-C277-teiden}



\begin{quote}
Before the new one: say the Japanese for to accumulate, then the Japanese for a snowstorm.
\end{quote}

\subsection*{You'll want to know: \ja{ていでん}}
\textbf{\ja{ていでん}} — \emph{teiden} — \textquotedblleft{}a power cut\textquotedblright{}.

\begin{quote}
\textbf{\ja{あらしで} \ja{ていでんに} \ja{なりました。}} — \emph{arashi de teiden ni narimashita} — \textquotedblleft{}The storm cut the power.\textquotedblright{}
\end{quote}

\begin{cousinweb}
\textbf{\ja{ていでん}} (\emph{teiden}) is Sino-Japanese (kango): two Chinese characters read with their on'yomi, readings borrowed from Chinese. They are \emph{tei}, \textquotedblleft{}halt\textquotedblright{} (Middle Chinese \emph{deng}) and \emph{den}, \textquotedblleft{}electricity\textquotedblright{} (\emph{denH}) (KANJIDIC2; Middle Chinese in Baxter's transcription).
\end{cousinweb}

\begin{grammarlens}[title={What you've built}]
One more word for telling what happened.
\end{grammarlens}

\subsection*{Your turn}
\begin{itemize}
\raggedright
  \item \emph{Say it:} \emph{teiden}
  \item \emph{Say it:} \emph{teiden}, once more
  \item \emph{Say it:} \emph{arashi de teiden ni narimashita} — The storm cut the power.
  \item \emph{Recall:} say the Japanese for to accumulate, then the Japanese for a snowstorm, then say \emph{teiden} again
\end{itemize}

\begin{tcolorbox}[breakable,colback=teal!4,colframe=teal!35!black,title={Before you move on}]
What does \ja{ていでん} mean? (\textquotedblleft{}A power cut\textquotedblright{}.) Say it once more.
\end{tcolorbox}
\section[\texorpdfstring{\ja{こごえる} (kogoeru)}{こごえる (kogoeru)}]{\ja{こごえる} (kogoeru) — to go numb with cold; to freeze}

\label{lesson:JA-C277-kogoeru}



\begin{quote}
Before the new one: say the Japanese for a snowstorm, then the Japanese for a power cut.
\end{quote}

\subsection*{You'll want to know: \ja{こごえる}}
\textbf{\ja{こごえる}} — \emph{kogoeru} — \textquotedblleft{}to go numb with cold; to freeze\textquotedblright{}.

\begin{quote}
\textbf{\ja{ゆきの} \ja{なかで} \ja{あしが} \ja{こごえました。}} — \emph{yuki no naka de ashi ga kogoemashita} — \textquotedblleft{}My feet went numb in the snow.\textquotedblright{}
\end{quote}

Polite \textbf{\ja{こごえます}}, \emph{kogoemasu}; an ichidan verb. Past (polite): \textbf{\ja{てが} \ja{こごえました。}}, \emph{te ga kogoemashita}.

\begin{cousinweb}
\textbf{\ja{こごえる}} (\emph{kogoeru}) is a native Japanese word (wago), inherited rather than borrowed from Chinese; its kanji was attached later as a kun'yomi (JMdict).
\end{cousinweb}

\begin{grammarlens}[title={What you've built}]
One more word for telling what happened.
\end{grammarlens}

\subsection*{Your turn}
\begin{itemize}
\raggedright
  \item \emph{Say it:} \emph{kogoeru}
  \item \emph{Say it:} \emph{kogoeru}, once more
  \item \emph{Say it:} \emph{yuki no naka de ashi ga kogoemashita} — My feet went numb in the snow.
  \item \emph{Recall:} say the Japanese for a snowstorm, then the Japanese for a power cut, then say \emph{kogoeru} again
\end{itemize}

\begin{tcolorbox}[breakable,colback=teal!4,colframe=teal!35!black,title={Before you move on}]
What does \ja{こごえる} mean? (\textquotedblleft{}To go numb with cold; to freeze\textquotedblright{}.) Say it once more.
\end{tcolorbox}
